% The abstract of the notes: overview, learning objectives, prerequisites.

\emph{Overview}
Coding agents now write much of the code that gets written, so who is a
literate program for?
We ask what harnesses you have, go through what research since 2023 says about
literate programming and generative AI, and look at how a coding agent can be
made to write literate programs.
We then give an agent the same task with three different instructions and judge
the results against the two goals of literate programming.
Finally we rate related approaches against the same goals: executable
demonstrations, instruction files for agents and specifications.

\emph{Learning objectives}
After this module, you can:
\begin{restatable}{lo}{AiLOReader}\label{AiLOReader}%
  Explain what prose in a program does for a human reader and what it does for
  a coding agent, and why these need not be the same thing.
\end{restatable}
\begin{restatable}{lo}{AiLOEvidence}\label{AiLOEvidence}%
  Judge a claim about AI and literate programming by the kind of evidence
  behind it: a proposal, a benchmark, or a study with people.
\end{restatable}
\begin{restatable}{lo}{AiLOInstruct}\label{AiLOInstruct}%
  Instruct a coding agent to write a literate program, and judge the result
  against the two goals of literate programming.
\end{restatable}
\begin{restatable}{lo}{AiLOCompare}\label{AiLOCompare}%
  Compare approaches related to literate programming by which of its goals they
  fulfil.
\end{restatable}

\emph{Prerequisites}
You should have worked through the lecture \emph{What is literate
programming?}, which supplies the two goals and noweb's tangling and weaving.
Nothing you are asked to do here needs a paid harness or a subscription.
